\section{Dynamic Programming from First Principles}

\begin{frame}{Calibration}
	\audience[Show of hands]{%
		\begin{enumerate}
			\item Have you implemented the convex hull trick or a Li Chao tree?
			\item Can you state the Bellman optimality equation from memory?
			\item Is Dijkstra's algorithm a dynamic program?
		\end{enumerate}
	}
\end{frame}

\begin{frame}{A motivating problem: minimum edit script}
	Given two versions of a file as token sequences $x \in \Sigma^m$ and $y \in \Sigma^n$,
	find a cheapest sequence of edits that turns $x$ into $y$.

	\begin{itemize}
		\item \textbf{delete} a token of $x$ \hfill cost $1$
		\item \textbf{insert} a token of $y$ \hfill cost $1$
		\item \textbf{substitute} a token \hfill cost $1$ (cost $0$ if the tokens already agree)
	\end{itemize}

	\seprule

	Goal: the minimum total cost $\mathrm{ED}(x,y)$, and an optimal script.
\end{frame}

\begin{frame}{State spaces}
	\begin{block}{Definition}
		A \alert{state space} is a weighted digraph $G=(V,E,\cost)$ with source $s$ and target $t$ such that
		solutions correspond to $s\to t$ paths and the objective is the sum of edge weights.
	\end{block}

	\begin{itemize}
		\item A state $v$ summarises a \emph{partial solution}.
		\item \alert{Sufficiency (Markov property):} the feasible completions of $v$, and
		their costs, depend only on $v$, not on the path that reached it.
	\end{itemize}
\end{frame}

\begin{frame}{The edit-distance state space}
	State $(i,j)$: ``$x_{1..i}$ has been transformed into $y_{1..j}$''.
	Source $(0,0)$, target $(m,n)$.
	\[
		\begin{array}{lll}
			(i-1,j)\to(i,j) & \text{delete } x_i & \cost=1\\
			(i,j-1)\to(i,j) & \text{insert } y_j & \cost=1\\
			(i-1,j-1)\to(i,j) & \text{match / substitute} & \cost=\indic{x_i\neq y_j}
		\end{array}
	\]
\end{frame}

\begin{frame}{The state space is a DAG}
	\begin{columns}[T,onlytextwidth]
		\column{0.58\textwidth}
		\centering
		\begin{tikzpicture}[x=1.15cm,y=0.95cm]
			\foreach \i in {0,...,4}{\foreach \j in {0,...,4}{
				\node[st] (n-\i-\j) at (\j,-\i) {};
			}}
			\foreach \i in {1,...,4}{\foreach \j in {0,...,4}{
				\pgfmathtruncatemacro{\ip}{\i-1}
				\draw[e] (n-\ip-\j) -- (n-\i-\j);
			}}
			\foreach \i in {0,...,4}{\foreach \j in {1,...,4}{
				\pgfmathtruncatemacro{\jp}{\j-1}
				\draw[e] (n-\i-\jp) -- (n-\i-\j);
			}}
			\foreach \i in {1,...,4}{\foreach \j in {1,...,4}{
				\pgfmathtruncatemacro{\ip}{\i-1}
				\pgfmathtruncatemacro{\jp}{\j-1}
				\draw[e] (n-\ip-\jp) -- (n-\i-\j);
			}}
			\draw[hl] (n-0-0) -- (n-1-0);
			\draw[hl] (n-1-0) -- (n-2-1);
			\draw[hl] (n-2-1) -- (n-3-2);
			\draw[hl] (n-3-2) -- (n-4-3);
			\draw[hl] (n-4-3) -- (n-4-4);
			\foreach \j/\c in {1/l,2/a,3/w,4/n}{\node[font=\scriptsize] at (\j,0.5) {$\c$};}
			\foreach \i/\c in {1/f,2/l,3/a,4/w}{\node[font=\scriptsize] at (-0.5,-\i) {$\c$};}
			\node[font=\scriptsize] at (-0.5,0.5) {$x\backslash y$};
		\end{tikzpicture}

		\column{0.4\textwidth}
		\begin{itemize}
			\item Every edge increases $i+j$: \alert{acyclic}.
			\item Highlighted: \texttt{flaw} $\to$ \texttt{lawn}, cost $2$.
			\item $(m+1)(n+1)$ states, $\le 3$ in-edges each.
		\end{itemize}
	\end{columns}
\end{frame}

\begin{frame}{Paths versus states}
	\audience[How many edit scripts?]{%
		For $m=n=20$, roughly how many $s\to t$ paths are there?\\[0.5em]
		(a)~$10^{6}$ \quad (b)~$10^{9}$ \quad (c)~$10^{12}$ \quad (d)~$>10^{14}$
	}
\end{frame}

\begin{frame}{Paths versus states: answer}
	The number of paths is the central Delannoy number $D(n,n)\sim c\,(3+2\sqrt2)^n$.
	\begin{center}
		$D(20,20) = 260{,}543{,}813{,}797{,}441 \approx 2.6\cdot 10^{14}$ paths
		\quad versus \quad $(20+1)^2=441$ states.
	\end{center}
	The state graph is a \alert{compressed representation} of the solution space.
\end{frame}

\begin{frame}{Dynamic programming on a DAG}
	\begin{theorem}
		Let $G=(V,E,\cost)$ be a DAG with source $s$, and $d(v)$ the minimum weight of an
		$s\to v$ path ($\infty$ if none). Then
		\[
			d(s)=0, \qquad d(v)=\min_{(u,v)\in E}\bigl\{\,d(u)+\cost(u,v)\,\bigr\}\quad (v\neq s).
		\]
	\end{theorem}
\end{frame}

\begin{frame}{Why it works}
	An $s\to v$ path is an $s\to u$ path followed by an edge $(u,v)$.
	\begin{itemize}
		\item \alert{Optimal substructure.} The choice of prefix does not constrain the last edge,
		so minimising each prefix loses nothing.
		\item \alert{Well-foundedness.} Acyclicity lets us induct along a topological order.
	\end{itemize}
\end{frame}

\begin{frame}{A strategy for dynamic programming problems}
	\begin{enumerate}
		\item \textbf{Decompose.} Choose a state space: what must be remembered about a
		partial solution so that the future is independent of the past?
		\item \textbf{Recurrence.} State the \alert{base case} and the \alert{inductive step}.
		\item \textbf{Evaluate bottom-up} in a topological order of the state DAG.
		\item \textbf{Reconstruct} the solution, if needed, from the minimising predecessor of each state.
	\end{enumerate}
\end{frame}

\begin{frame}{Complexity is read off the DAG}
	\[
		T \;=\; \sum_{v\in V} \bigl(\text{work to evaluate state } v\bigr) \;=\; |V|\times(\text{cost per state}) \;\approx\; |E|.
	\]
\end{frame}

\begin{frame}{Steps 1--2: the recurrence for edit distance}
	Let $f(i,j)$ be the minimum cost of transforming $x_{1..i}$ into $y_{1..j}$.

	\begin{align*}
		\text{base:}\quad & f(i,0)=i, \qquad f(0,j)=j\\[0.3em]
		\text{step:}\quad & f(i,j)=\min\begin{cases}
			f(i-1,j)+1 & \text{delete } x_i\\
			f(i,j-1)+1 & \text{insert } y_j\\
			f(i-1,j-1)+\indic{x_i\neq y_j} & \text{match / substitute}
		\end{cases}
	\end{align*}
	Answer: $\mathrm{ED}(x,y)=f(m,n)$.
\end{frame}

\begin{frame}{Step 3: bottom-up evaluation}
	\audience[Fill in the last two rows]{%
		\centering
		\renewcommand{\arraystretch}{1.15}
		\begin{tabular}{@{}c|ccccc@{}}
			& $\varepsilon$ & \texttt{l} & \texttt{a} & \texttt{w} & \texttt{n}\\
			\midrule
			$\varepsilon$ & 0 & 1 & 2 & 3 & 4\\
			\texttt{f} & 1 & 1 & 2 & 3 & 4\\
			\texttt{l} & 2 & 1 & 2 & 3 & 4\\
			\texttt{a} & ? & ? & ? & ? & ?\\
			\texttt{w} & ? & ? & ? & ? & ?\\
		\end{tabular}
	}
\end{frame}

\begin{frame}{Step 3: bottom-up evaluation, filled in}
	\centering
	\renewcommand{\arraystretch}{1.15}
	\begin{tabular}{@{}c|ccccc@{}}
		& $\varepsilon$ & \texttt{l} & \texttt{a} & \texttt{w} & \texttt{n}\\
		\midrule
		$\varepsilon$ & 0 & 1 & 2 & 3 & 4\\
		\texttt{f} & 1 & 1 & 2 & 3 & 4\\
		\texttt{l} & 2 & 1 & 2 & 3 & 4\\
		\texttt{a} & 3 & 2 & \alert{1} & 2 & 3\\
		\texttt{w} & 4 & 3 & 2 & \alert{1} & \alert{2}\\
	\end{tabular}

	\vspace{1em}
	$\mathrm{ED}(\texttt{flaw},\texttt{lawn})=f(4,4)=2$.
\end{frame}

\begin{frame}{Step 4: reconstruction}
	Walk back from $(m,n)$ along minimising predecessors:
	\[
		(4,4)\xleftarrow{\text{insert n}}(4,3)\xleftarrow{\text{match w}}(3,2)\xleftarrow{\text{match a}}(2,1)\xleftarrow{\text{match l}}(1,0)\xleftarrow{\text{delete f}}(0,0).
	\]
	\begin{center}
		\ttfamily
		\begin{tabular}{@{}cccccc@{}}
			f & l & a & w & - \\
			- & l & a & w & n \\
			\textrm{del} & = & = & = & \textrm{ins}
		\end{tabular}
	\end{center}
\end{frame}

\begin{frame}{Memory}
	\begin{itemize}
		\item The value table costs $\Theta(mn)$ memory; the argmin can be recomputed from it.
		\item Rolling two rows: $\Theta(\min(m,n))$ memory, but no reconstruction.
		\item \alert{Hirschberg's divide-and-conquer}: reconstruction in linear space at twice the time.
	\end{itemize}
\end{frame}

\begin{frame}{Complexity and its limits}
	$|V|=(m+1)(n+1)$, at most $3$ predecessors each: $\Theta(mn)$ time.

	\begin{itemize}
		\item \textbf{Output-sensitive.} If $\mathrm{ED}\le d$, only the band $|i-j|\le d$ matters:
		$O(d\cdot\min(m,n))$ (Ukkonen); Myers' $O(nd)$ diff is the engine behind \texttt{git diff}.
		\item \textbf{Lower bound.} No $O(n^{2-\varepsilon})$ algorithm exists unless SETH
		fails (Backurs--Indyk, STOC 2015).
	\end{itemize}
\end{frame}

\begin{frame}{Is the state space really sufficient?}
	\audience[Adding transpositions]{%
		Allow \textbf{transpositions} of adjacent tokens ($ab\to ba$) at cost $1$.
		\begin{enumerate}
			\item Which edge do you add to the DAG, and when?
			\item Does the complexity change?
			\item Is the result exactly the Damerau--Levenshtein distance?
		\end{enumerate}
	}
\end{frame}

\begin{frame}{Transpositions: the edge}
	\begin{enumerate}
		\item Add $(i-2,j-2)\to(i,j)$ with cost $1$ when $x_i=y_{j-1}$ and $x_{i-1}=y_j$.
		\item Still $\Theta(mn)$: one more $O(1)$ in-edge per state.
	\end{enumerate}
\end{frame}

\begin{frame}{Transpositions: not Damerau--Levenshtein}
	\begin{enumerate}
		\setcounter{enumi}{2}
		\item \alert{No.} This computes the \emph{optimal string alignment} distance, in which no substring is edited twice.
	\end{enumerate}

	\vspace{0.5em}
	For $x=\texttt{CA}$, $y=\texttt{ABC}$ it gives $3$.
	The true distance is $2$: \texttt{CA}$\to$\texttt{AC}$\to$\texttt{ABC}.

	\seprule

	The state space is \alert{not sufficient}. The exact algorithm needs edges reaching back to the
	last occurrence of a symbol (see backup).
\end{frame}

\begin{frame}{Part 1 in one slide}
	What made this a dynamic program?
	\begin{enumerate}
		\item \textbf{A sufficient state}, giving a graph that compresses exponentially many solutions.
		\item \textbf{Acyclicity}, giving a well-founded order of evaluation.
		\item \textbf{Additive (or otherwise decomposable) objective}, giving optimal substructure.
	\end{enumerate}

	\seprule

	\textbf{Design is in step 1.} Steps 2--4 and the complexity follow mechanically.
\end{frame}
